\documentclass[12pt, twoside, openright]{report} % Fuente a 12pt, formato doble página y chapter a la derecha
\raggedbottom % No ajustar el contenido con un salto de página

% MÁRGENES: 2,5 cm sup. e inf.; 3 cm izdo. y dcho.
\usepackage[
a4paper,
vmargin=2.5cm,
hmargin=3cm
]{geometry}

% INTERLINEADO: Estrecho (6 ptos./interlineado 1,15) o Moderado (6 ptos./interlineado 1,5)
\renewcommand{\baselinestretch}{1.15}
\parskip=6pt

% DEFINICIÓN DE COLORES para portada y listados de código
\usepackage[table]{xcolor}
\definecolor{azulUC3M}{RGB}{0,0,102}
\definecolor{gray97}{gray}{.97}
\definecolor{gray75}{gray}{.75}
\definecolor{gray45}{gray}{.45}

% Soporte para GENERAR PDF/A
\usepackage{etoolbox}
\makeatletter
\@ifl@t@r\fmtversion{2021-06-01}%
 {\AddToHook{package/after/xmpincl}
   {\patchcmd\mcs@xmpincl@patchFile{\if\par}{\ifx\par}{}{\fail}}}{}
\makeatother
\usepackage[a-1b]{pdfx}

% ENLACES
\usepackage{hyperref}
\hypersetup{colorlinks=true,
  linkcolor=black, % enlaces a partes del documento (p.e. índice) en color negro
  urlcolor=blue} % enlaces a recursos fuera del documento en azul

% Añadir pdfs como partes del documento
\usepackage{pdfpages}

% Quitar la indentación de principio de los párrafos
\setlength{\parindent}{0em}

% EXPRESIONES MATEMÁTICAS
\usepackage{amsmath,amssymb,amsfonts,amsthm}

\usepackage{txfonts} 
\usepackage[T1]{fontenc}
\usepackage[utf8]{inputenc}

% Insertar gráficas y fotos
\usepackage{tikz}
\usetikzlibrary{cd}
\usepackage{pgfplots}

\usepackage[spanish, es-tabla]{babel} 
\usepackage[babel, spanish=spanish]{csquotes}
\AtBeginEnvironment{quote}{\small}

% diseño de PIE DE PÁGINA
\usepackage{fancyhdr}
\pagestyle{fancy}
\fancyhf{}
\renewcommand{\headrulewidth}{0pt}
\fancyfoot[LE,RO]{\thepage}
\fancypagestyle{plain}{\pagestyle{fancy}}

% DISEÑO DE LOS TÍTULOS de las partes del trabajo (capítulos y epígrafes o subcapítulos)
\usepackage{titlesec}
\usepackage{titletoc}
\titleformat{\chapter}[block]
{\large\bfseries\filcenter}
{\thechapter.}
{5pt}
{\MakeUppercase}
{}
\titlespacing{\chapter}{0pt}{0pt}{*3}
\titlecontents{chapter}
[0pt]                                               
{}
{\contentsmargin{0pt}\thecontentslabel.\enspace\uppercase}
{\contentsmargin{0pt}\uppercase}                        
{\titlerule*[.7pc]{.}\contentspage}                 

\titleformat{\section}
{\bfseries}
{\thesection.}
{5pt}
{}
\titlecontents{section}
[5pt]                                               
{}
{\contentsmargin{0pt}\thecontentslabel.\enspace}
{\contentsmargin{0pt}}
{\titlerule*[.7pc]{.}\contentspage}

\titleformat{\subsection}
{\normalsize\bfseries}
{\thesubsection.}
{5pt}
{}
\titlecontents{subsection}
[10pt]                                               
{}
{\contentsmargin{0pt}                          
  \thecontentslabel.\enspace}
{\contentsmargin{0pt}}                        
{\titlerule*[.7pc]{.}\contentspage}  

\usepackage{adjustbox}

% DISEÑO DE TABLAS.
\usepackage{multirow} % permite combinar celdas 
\usepackage{caption} % para personalizar el título de tablas y figuras
\usepackage{floatrow} % utilizamos este paquete y sus macros \ttabbox y \ffigbox para alinear los nombres de tablas y figuras de acuerdo con el estilo definido. Para su uso ver archivo de ejemplo 
\usepackage{array} % con este paquete podemos definir en la siguiente línea un nuevo tipo de columna para tablas: ancho personalizado y contenido centrado
\newcolumntype{P}[1]{>{\centering\arraybackslash}p{#1}}
\DeclareCaptionFormat{upper}{#1#2\uppercase{#3}\par}

% Diseño de tabla para ingeniería
\captionsetup[table]{
  format=hang,
  name=Tabla,
  justification=centering,
  labelsep=colon,
  width=.75\linewidth,
  labelfont=small,
  font=small,
}

% DISEÑO DE FIGURAS.
\usepackage{graphicx}
\graphicspath{{img/}} %ruta a la carpeta de imágenes

% Diseño de figuras para ingeniería
\captionsetup[figure]{
  format=hang,
  name=Fig.,
  singlelinecheck=off,
  labelsep=colon,
  labelfont=small,
  font=small    
}

% NOTAS A PIE DE PÁGINA
\usepackage{chngcntr} % Para numeración continua de las notas al pie
\counterwithout{footnote}{chapter}

% LISTADOS DE CÓDIGO
% soporte y estilo para listados de código. Más información en https://es.wikibooks.org/wiki/Manual_de_LaTeX/Listados_de_código/Listados_con_listings
\usepackage{listings}

% definimos un estilo de listings
\lstdefinestyle{estilo}{ frame=Ltb,
  framerule=0pt,
  aboveskip=0.5cm,
  framextopmargin=3pt,
  framexbottommargin=3pt,
  framexleftmargin=0.4cm,
  framesep=0pt,
  rulesep=.4pt,
  backgroundcolor=\color{gray97},
  rulesepcolor=\color{black},
  %
  basicstyle=\ttfamily\footnotesize,
  keywordstyle=\bfseries,
  stringstyle=\ttfamily,
  showstringspaces = false,
  commentstyle=\color{gray45},     
  %
  numbers=left,
  numbersep=15pt,
  numberstyle=\tiny,
  numberfirstline = false,
  breaklines=true,
  xleftmargin=\parindent
}

\captionsetup[lstlisting]{font=small, labelsep=period}
% fijamos el estilo a utilizar 
\lstset{style=estilo}
\renewcommand{\lstlistingname}{\uppercase{Código}}

\pgfplotsset{compat=1.17} 
%-------------
% DOCUMENTO
%-------------

\begin{document}
\pagenumbering{roman} % Se utilizan cifras romanas en la numeración de las páginas previas al cuerpo del trabajo

%----------
% PORTADA
%---------- 
\begin{titlepage}
  \begin{sffamily}
    \color{azulUC3M}
    \begin{center}
      \begin{figure}[H] % Incluimos el logotipo de la Universidad
        \makebox[\textwidth][c]{\includegraphics[width=16cm]{Portada_Logo.png}}
      \end{figure}
      \vspace{2.5cm}
      \begin{Large}
        Grado en Ingeniería Informática\\
        2021-2022\\
        \vspace{2cm}
        \textsl{Apuntes}\\
        \bigskip
      \end{Large}
      {\Huge Informática Gráfica}\\
      \vspace*{0.5cm}
      \rule{10.5cm}{0.1mm}\\
      \vspace*{0.9cm}
      {\LARGE Jorge Rodríguez Fraile\footnote{\href{mailto:100405951@alumnos.uc3m.es}{Universidad: 100405951@alumnos.uc3m.es}  |  \href{mailto:jrf1616@gmail.com}{Personal: jrf1616@gmail.com}}}\\
      \vspace*{1cm}
    \end{center}
    \vfill
    \color{black}
    \includegraphics[width=4.2cm]{img/creativecommons.png}\\
    Esta obra se encuentra sujeta a la licencia Creative Commons\\ \textbf{Reconocimiento - No Comercial - Sin Obra Derivada}
  \end{sffamily}
\end{titlepage}

%----------
% ÍNDICES
%----------

%--
% Índice general
%-
\tableofcontents
\thispagestyle{fancy}

%--
% Índice de figuras. Si no se incluyen, comenta las líneas siguientes
%-
\listoffigures
\thispagestyle{fancy}

%--
% Índice de tablas. Si no se incluyen, comenta las líneas siguientes
%-
\listoftables
\thispagestyle{fancy}

%----------
% TRABAJO
%----------

\pagenumbering{arabic} % numeración con números arábigos para el resto de la publicación

%----------
% COMENZAR A ESCRIBIR AQUÍ
%----------

\chapter{Información}\label{ch:informacion}
\section{Profesores}\label{sec:profesores}
\begin{quote}
    Coordinador: Antonio Berlanga

	  Magistral: Yago Sáez, yago.saez@uc3m.es, 2.1.C.13

    Magistral: José María Valls

    Prácticas: David Quintana
\end{quote}

\section{Sistema de evaluación}\label{sec:sistema-de-evaluación}
\begin{itemize}
  \item 50 \% Teoría
  \begin{itemize}
    \item 10 \% Test 1. 4 temas
    \item 10 \% Test 2. 4 temas
    \item 30 \% Examen Final. Todos los contenidos teóricos de la asignatura \\ Nota mínima: 3
  \end{itemize}
  \item 50 \% Práctica. La 1, 2 y 3 son por parejas y la última en grupo de 3
  \begin{itemize}
    \item 12,5 \% Práctica 1. Modelado 3D
    \item 10 \% Práctica 2. Fractales
    \item 12,5 \% Práctica 3. Paisajes y terreno
    \item 15 \% Práctica Final. Modelado 3D, Iluminación y sombreado y Animación
  \end{itemize}
\end{itemize}

\chapter{Introducción a la Informática Gráfica}\label{ch:introduccion-a-la-informática-gráfica}
\section{Que es una Imagen}\label{sec:que-es-una-imagen}
\begin{itemize}
	\item Una vista en un monitor.
	\item Un fichero en un cámara.
	\item Números de una RAM
\end{itemize}
\textbf{Definición}
\begin{itemize}
	\item Una distribución 2D de intensidad o color.
	\item Una función definida en un plano bidimensional.
	\item ! No hay mención a píxeles!
\end{itemize}
\textbf{Necesidades}
\begin{itemize}
	\item Representar la imagen (codificarla en números)
	\item Mostrar la imagen (aplicar transformaciones, intensidades de corriente, cantidades de tinta).
\end{itemize}

\section{Áreas de Conocimiento}\label{sec:áreas-de-conocimiento}
\textbf{Procesado de imágenes}, se reciben solamente imágenes, se hacen modificaciones y análisis de imágenes.
\begin{itemize}
	\item Análisis de escenas.
	\item Reconstrucción de modelos en 2D o 3D de las imágenes.
	\item Compresión de imágenes, mejora de nitidez, posterizado (convertir imágenes en dibujos o pósteres)\ldots
\end{itemize}

\textbf{Visión Artificial}, dadas imágenes y datos sobre estas se obtiene información más completa.
\begin{itemize}
	\item Construcción de sistemas que obtengan información automáticamente a partir de las imágenes.
	\item Clasificación de imágenes en base a conocimiento previo o información estadística.
\end{itemize}

\textbf{Informática Gráfica}, síntesis de objetos reales o imaginarios a partir de modelos basados en computación.

\section{Aplicaciones}\label{sec:aplicaciones}

\subsection{Diseño Asistido por Computador (CAD)}\label{subsec:diseño-asistido-por-computador-(cad)}
Herramientas gráficas para diseñar prototipos y evaluarlos antes de construirlos.

\textbf{Áreas importantes:} Diseño industrial, Arquitectura, Circuitos electrónicos y electrónicos.

\textbf{Técnicas:} Diseño basado en primitivas constructivas y Superficies curvas.

\textbf{Posibilidades:} Realidad virtual, Presentación realista, Análisis del diseño y Conexión con el sistema de fabricación.

\subsection{Gráficos de presentación}\label{subsec:gráficos-de-presentación}
Uso de los gráficos para producción de ilustraciones de soporte a informes y trabajos.

\textbf{Áreas importantes:} Economía, Estadística, Matemáticas y Administración y gestión.

\textbf{Técnicas:} Gráficos de líneas, Gráficos de barra, Gráficos de tarta y Superficies 3D\@.

\subsection{Creaciones artisticas}\label{subsec:creaciones-artisticas}
Producción de imágenes con un fin artístico o comercial.

\textbf{Áreas importantes:} Diseño de logotipos, Bellas artes y Animaciones publicitarias.

\textbf{Técnicas:} Diseño vectorial, Soporte a la animación, Tratamiento de imagen y Rendering.

\subsection{Entretenimiento}\label{subsec:entretenimiento}
Producción de videos y videojuegos con efectos realistas.

\textbf{Áreas importantes:} Cine (películas y animación), Televisión (cortinillas, cabeceras, AR) y Videojuegos.

\textbf{Técnicas:} Animación, Visualización realista, Efectos especiales e Interactividad.

Se utilizan API y SDK: Truevision 3D, XNA, OGRE, IrrLicht, CryEngine 3, Unreal Engine y Unity.

\subsection{Simulación y entrenamiento}\label{subsec:simulación-y-entrenamiento}
Desarrollo de sistemas para simular tareas realistas y procesos.

\textbf{Áreas importantes:} Simulación de conducción, Simulación de procesos industriales, Entrenamiento y Educación.

\textbf{Técnicas:} Tiempo real, Interactividad y Realidad virtual.

\subsection{Visualización cientifica}\label{subsec:visualización-cientifica}
Visualización en dominios específicos y de grandes cantidades de datos.

\textbf{Áreas importantes:} Medicina (resonancias, tomografías), Ingeniería, Física (campos, dinámica de fuerza), Química (interacción moléculas), Matemática y Topológico (terrenos y corrientes).

\textbf{Técnicas:} Codificación por color, Curvas de nivel, Visualización de volúmenes.

\section{Historia}\label{sec:historia}
\subsection{Años 50 (Comienzos)}
Whirlwind y SA\@.E.

Se empezó utilizando la informática gráfica para detectar aviones o submarinos, con radares, en estos tiempos se utilizaban ordenadores del tamaño de ordenadores. En esta visualización se podía marcar puntos.

\subsection{Años 60}\label{subsec:años-60}
\begin{itemize}
	\item Sketchpad: el primer programa gráfico interactivo.
	\item Russell (MIT) desarrolla Spacewar en un PDP-1.
	\item El primer juego ''Tennis for two'' (Pong) se jugaba en un osciloscopio.
	\item Primer algoritmo de superficies ocultas (Catmull)
	\item Realismo mediante sombreado de superficies con color.
\end{itemize}

\subsection{Años 70}\label{subsec:años-70}
\begin{itemize}
	\item Suavizado de superficies poligonales (Gouraud)
	\item Comercialización del microprocesador.
	\item Fundación de Atari.
	\item Primeros intentos de informática gráfica en el cine.
	\item Diseños de la tetera (Newell, Utah, Univ.)
	\item Introducción de texturas y Z-buffer.
	\item Suavizado de superficies poligonales (Phong) que se utiliza en el 3D
	\item Fundación de Apple y Microsoft.
	\item Lucasfilm crea la división de gráficos por computador.
	\item Westworld: primera película en emplear gráficos por ordenador.
	\item La Guerra de las Galaxias.
\end{itemize}

\subsection{Anos 80}\label{subsec:anos-80}
\begin{itemize}
	\item Popularización de SIGGRAPH como evento de referencia en el área.
	\item Publicación sobre el Raytracing (Bell Labs)
	\item Construcción del primer motor de rendering (REYES), precursor de Renderman (Carpenter)
	\item Ecuación de rendering (Kajiya)
	\item Película de TRON de Disney (Lisberger y Kushner)
	\item Venta masiva de terminales gráficas: IBM, Tektronix\ldots
	\item GKS se constituye como estándar ISO y ANSI de construcción de librerías gráficas.
	\item IBM crea el ordenador personal (PC)
	\item Introducción de la radiosidad (Goral, Torrance, Cohen)
	\item Cornell Box: prueba para comprobar la adecuación de un sistema de renderizado a una fotografía real.
\end{itemize}

\section{Campos de aplicación}\label{sec:campos-de-aplicación}
\begin{itemize}
	\item \textbf{Física, matemáticas y ciencias naturales:} Simulación, Modelado y Análisis numérico.
	\item \textbf{Ingeniería:} Sistemas software y hardware, e Infraestructura e integración de entonos.
	\item \textbf{Arte y psicología:} Percepción y Diseño y composición.
\end{itemize}

\section{Disciplinas relacionadas}\label{sec:disciplinas-relacionadas}
\begin{figure}[H]
	\ffigbox[\FBwidth]
	{\caption{Disciplinas relacionadas}}
	{\def\svgwidth{.8\textwidth}
		\input{./img/disciplinas-relacionadas.eps_tex}}
\end{figure}

\section{Procesado de Imagen 2D}\label{sec:procesado-de-imagen-2d}
La secuencia de pasos por los que pasa el procesado de imagen en 2D puede representarse como un pipeline (no todos los pasos tienen que estar presentes)
\begin{figure}[H]
	\ffigbox[\FBwidth]
	{\caption{Pipeline Imagenes 2D}}
	{\def\svgwidth{.8\textwidth}
		\input{./img/pipeline2D.eps_tex}}
\end{figure}

\begin{enumerate}
	\item \textbf{Adquisición:} Muchas técnicas, divisibles en dos categorías: Síntesis y Captura.
	      \begin{itemize}
		      \item \textbf{Síntesis:} Imágenes creadas por un ordenador. \\ Definición geométrica de los objetos (Renderman, Maya, PovRay). \\ Imágenes "pintadas"(Photoshop, Illustrator, Fractal Painter).
		      \item \textbf{Captura:} Imágenes que provienen del mundo real. \\ Capturadas o digitalizadas por hardware específico.
	      \end{itemize}
	\item \textbf{Preprocesado:} Modificación de formas, tamaño y propiedades de color. \\ Técnicas: Ajuste de color y niveles, Recorte, Escalado, Desenfoque, Mejora de bordes y Filtrado.
	\item \textbf{Mapeado:} Varias imágenes se combinan con transformaciones. \\ Transformaciones: Rotaciones, Escalados y Deformaciones. \\ Composición: Efectos de transparencia / translucencia.
	\item \textbf{Postprocesado:} Se utiliza para aplicar efectos globales a toda o parte de la imagen. \\ Efectos artísticos: Posterizado, Envejecimiento, Desenfoque  o Texturizado. \\ Efectos técnicos: Mejora de contraste y Variaciones de color.
	\item \textbf{Salida:} El dispositivo de salida puede afectar a su visualización.
	      \begin{itemize}
		      \item El mapa de color de una impresora puede <"falsear>" o acentuar colores.
		      \item Se pueden mapear los colores entre la pantalla y la impresora.
	      \end{itemize}
	      Dispositivos: Monitores, Impresoras, Discos y Mapas de texturas.
\end{enumerate}

\section{Procesado Gráfico 3D}\label{sec:procesado-gráfico-3d}
Se encarga de generar las imágenes que se pueden ver en un dispositivo.
\begin{itemize}
	\item Vista del sistema operativo.
	\item Imágenes (frames) de un videojuego.
	\item Pantalla de un teléfono.
	\item Imágenes de una película en el cine.
\end{itemize}
\subsection{Pipeline 3D}\label{subsec:pipeline-3d}
Varios módulos van realizando diferentes tareas para la generación (render) de la imagen.
\begin{itemize}
	\item \textbf{Transformación}
	\item \textbf{Iluminación}
	\item \textbf{Texturado}
	\item \textbf{Shaders:} Unidades de procesado que se compilan de forma independiente, modificando localmente las propiedades de los objetos de la escena. \\ Vértices: Transformaciones de coordenadas y Colores y texturas. \\ Geométricos: Generación de formas (primitivas). \\ Fragmentos: Texturizado y sombreado a nivel de pixel.
\end{itemize}
Diagrama de flujo
\begin{itemize}
	\item Entrada de los datos de la geometría de objetos y atributos de sus materiales.
	\item Definición de las condiciones de iluminación.
	\item Salida de las imágenes.
\end{itemize}
Historia
\begin{itemize}
	\item Datos de entrada > Software de Render > Frame Buffer.
	\item Datos de entrada > Transformaciones e Iluminación > Transformaciones Raster > Procesado Pixel > Frame Buffer.
	\item Datos de entrada > Vertex Shading > Geometry Shading > Transformaciones Raster > Pixel Shading > Frame Buffer.
\end{itemize}

Entran vértices y salen píxeles.

\subsection{Framebuffer}\label{subsec:framebuffer}
Memoria local para las imágenes que se mapean en un dispositivo. El hardware de video convierte el contenido en una señal para el dispositivo de salida.

\textbf{Framebuffer sencillo:}
\begin{itemize}
	\item Dividido en planos.
	\item Cada plano representa 1 bit de valor del color del píxel.
	\item El plano tiene el tamaño del dispositivo (ejem. 1920x1080).
	\item Para RGB, profundidad de 24 bits.
\end{itemize}

\textbf{Framebuffer complejo}
\begin{itemize}
	\item Doblar el número de planos (double buffering).
	\item Z-buffer para profundidad.
	\item Visión 3D, un conjunto de buffers para cada ojo.
	\item A-buffer para transparencia y sombras.
\end{itemize}


Z-buffer: es importante la precisión utilizada. Pueden aparecer problemas de "z-fighting", que dos superficies que ocupen el mismo espacio y se vea como una se funde con la otra.

El framebuffer necesita memorias muy rápidas
\pagebreak

\subsection{Flujo de aplicación}\label{subsec:flujo-de-aplicación}
\begin{enumerate}
	\item Software
	      \begin{enumerate}
		      \item \textbf{Aplicación 3D:} CAD, Animaciones, Simulador o Juego 3D\@.
		      \item \textbf{API:} Direct3D u OpenGL\@.
		      \item \textbf{Drivers:} Preparación de datos, Mapeo hardware.
	      \end{enumerate}
	\item Hardware
	      \begin{enumerate}
		      \item \textbf{Motor gráfico:} Transformaciones geométricas, Ilustración, Texturizado y Rasterización.
		      \item \textbf{Memoria local:} Framebuffer, z-buffer y Stereo buffer.
	      \end{enumerate}
\end{enumerate}

\subsection{Motor de Render}\label{subsec:motor-de-render}
Software encargado de generar y visualizar información gráfica Abstrae los mecanismos de creación de gráficos.

Hay una gran variedad de motores gráficos en el mercado, de distintos tipos opensource, freeware y comerciales.

\textbf{Técnicas:} Sombreado, Mapeo de texturas, Bump-mapping, PBR, Efectos atmosféricos, Sombras, Reflexiones, Refracciones, Transparencias, Translucencias, Difracción, Iluminación global, Cáusticas, Profundidad de campo, Motion blur y NPR (cell-shading)
\chapter{Percepción Visual}
La percepción visual se divide en:
\begin{itemize}
	\item \textbf{Percepción física:} Cuando los fotones llegan a nuestro nervio óptico y se transforman en impulsos electromagnéticos que estimulan el cerebro.
	\item \textbf{Procesado e interpretación:} Se transforman los estímulos/impulsos en imagen.
\end{itemize}

Por las características físicas, los humanos no podemos verlo todo.

El sistema de interpretación puede construir imágenes a partir de información incompleta.

El procesamiento compensa el Movimiento y Variaciones en la luminosidad, pero el contexto es fundamental. Las expectativas resuelven ambigüedades.

\textbf{Movimiento:} Mínimo de 16 fotogramas/segundo para percibirlo. Al recibir las imágenes tan rápido sentimos que se mueven.
\begin{itemize}
	\item Cine mudo 16-18 Hz, Cine 24 Hz y TV 25-29 Hz.
\end{itemize}

\section{Ilusiones Ópticas}
El procesamiento visual puede crear ilusiones ópticas.
\begin{itemize}
	\item Nombres de colores de diferente color.
	\item Tamaño de objetos iguales, pero en una posición diferente, escaleras, flechas o platos.
	\item Objetos que se mueven o cambian al mirarlos.
\end{itemize}

\section{Percepción de la Luminosidad}
\textbf{Luminosidad:} Cantidad de luz que recibimos (subjetiva) y depende de:
\begin{itemize}
	\item \textbf{Luminancia:} Cantidad de luz emitida por un objeto.
	\item \textbf{Contraste:} Relación entre la luminancia de un objeto y la luminancia de su entorno.
\end{itemize}

\section{Percepción de Tamaño y Profundidad}
\textbf{Ángulo de visión:} Ángulo definido por los bordes del objeto percibido y el ojo. Los objetos familiares se perciben a tamaño constante.

\textbf{Agudeza visual:} Capacidad para percibir detalles. El límite está en un ángulo visual de 1/60 grados.

\textbf{Pistas visuales:} Estas ayudas permiten tener percepción del tamaño y profundidad.
\begin{itemize}
	\item Tamaño y posición relativa de los objetos.
	\item Más/menos definidos parecen estar más cerca o más lejos.
\end{itemize}

Percepción de tamaño con el mismo ángulo visual.

Tamaños iguales percibimos como diferentes - diferente tamaño relativo.

Tamaño iguales percibimos como diferentes - diferente profundidad.

Tamaño diferentes percibidos como iguales - diferente profundidad.

\section{Características de la Visión}
Percepción del tamaño y la profundidad
\begin{itemize}
	\item \textbf{Ángulo de visión:} Cantidad que un objeto ocupa el área visual. Permite determinar tamaño y distancia.
	\item \textbf{Agudeza visual:} Capacidad para percibir detalles.
	\item \textbf{Pistas:} Ayudas visuales que permiten tener percepción del tamaño y la profundidad.
\end{itemize}

Pistas visuales
\begin{itemize}
	\item \textbf{Contraste, claridad y brillo:} Objetos más definidos parecen más cercanos.
	\item \textbf{Sombras:} Pistas adicionales de la posición relativa.
	\item \textbf{Texturas:} Menos definidas cuando más alejados del observador.
\end{itemize}

\chapter{Proyecciones, Objetos 3D y GSD}
\section{Proyecciones}
\textbf{Proyección:} Transformación matemática que convierte un conjunto de puntos n-dimensional en un conjunto de puntos k-dimensional, siendo $k \leq n$. Reduce la dimensionalidad, normalmente 3D a 2D.

La proyección se define por unas \textbf{líneas de proyección} (proyectores) que, dirigidos hacia el \textbf{centro de proyección}, pasan a través del objeto e intersecan el \textbf{plano de proyección/superficie de proyección} para formar la proyección del objeto en el mismo.

Proyecciones Geométricas Planas
\begin{itemize}
    \item Paralela
    \begin{itemize}
        \item Ortográfica
        \begin{itemize}
            \item Elevaciones, las proyecciones ortográficas en tres planos.
            \item Axonométrica, el objeto es rotado respecto a uno o más ejes del plano de proyección.
            \begin{itemize}
                \item Isométrica
            \end{itemize}
        \end{itemize}
        \item Oblicua
        \begin{itemize}
            \item Cabinet, eje de profundidad (z) forma 45 grados con el eje x. Profundidad escalada (1/2, 2/3, etc.).
            \item Caballera, eje de profundidad (z) forma 45 grados con el eje x. Profundidad real.
        \end{itemize}
    \end{itemize}
    \item perspectiva
    \begin{itemize}
        \item 1 punto
        \item 2 puntos
        \item 3 puntos
    \end{itemize}
\end{itemize}

Métodos de generar vistas de objetos:
\begin{itemize}
    \item \textbf{Proyección paralela:} El centro de proyección en el infinito, los proyectores son paralelos hacia el centro.
    \begin{itemize}
        \item Mantiene la forma y dimensión de los objetos.
        \item Representación no realista de los objetos, se usan para el diseño.
        \item Conservan las propiedades relativas.
    \end{itemize}
    \begin{figure}[H]
        \ffigbox[\FBwidth]
        {\caption{Proyección perspectiva}}
        {\def\svgwidth{.5\textwidth}
            \input{./img/perspectiva.eps_tex}}
    \end{figure}
    \item \textbf{Proyección perspectiva:} Hay un centro de proyección cercano.
    \begin{itemize}
        \item Altera la forma y dimensión de los objetos.
        \item Representación realista de los objetos, simula la visión humana.
        \item No conserva las proporciones
    \end{itemize}
    \vspace{-5cm}
    \begin{figure}[H]
        \ffigbox[\FBwidth]
        {\caption{Proyección paralela}}
        {\def\svgwidth{.5\textwidth}
            \input{./img/paralela.eps_tex}}
    \end{figure}
\end{itemize}

\subsection{Visión Estereográfica}
Se obtiene realizando dos proyecciones en perspectiva, ambas imágenes se combinan para dar una tridimensional. Es el caso de los ojos.

\subsection{Proyecciones no planas}
\begin{itemize}
    \item Proyección esférica: Imagen proyectada sobre una esfera.
    \item Proyección fisheye (ojo de pez, media esfera): Proyección esférica con un ángulo de 180 grados (estándar) o 360 grados (ultra). La proyección sobre un círculo.
    \item Proyección ultra wide angle: Fisheye en el que la proyección es un rectángulo.
    \item Proyección Omnimax: Fisheye con un ángulo de visión reducido en el eje vertical.
    \item Proyección cilíndrica: Imagen proyectada sobre un cilindro.
    \item Proyección panorámica: Tipo de proyección cilíndrica capaz de ver ángulos superiores a 180 grados.
\end{itemize}

\section{Representación de Objetos 3D}
Definir estructuras de datos 3D capaces de representar las propiedades geométricas y físicas.

Debemos ser capaces de representar el interior, el exterior y el comportamiento de la superficie.

Aproximaciones al problema:
\begin{itemize}
    \item Definición de la superficie exterior: Modelos de superficie, solo interesa el contorno (exterior).
    \item Definición del espacio ocupado por el objeto: Modelado Sólido, para saber exterior, interior y propiedad de los objetos.
\end{itemize}

Requisitos:
\begin{itemize}
    \item Facilidad en la generación de objetos.
    \item Calidad en el aspecto final de la representación, según las restricciones.
    \item Recursos computacionales de la aplicación.
\end{itemize}

Representación de Objetos 3D
\begin{itemize}
    \item Modelos de Superficie
    \begin{itemize}
        \item Poligonales
        \item Matemáticos
        \begin{itemize}
            \item Superficies Implícitas
            \item Superficies Paramétricas
        \end{itemize}
    \end{itemize}
    \item Modelo Sólido
    \begin{itemize}
        \item Modelos de Barrido
        \item Geometría Sólida Constructiva (GSC)
    \end{itemize}
    \item Modelos de Partición Espacial
\end{itemize}

\subsection{Modelos de Superficie}
Se emplea cuándo
\begin{itemize}
    \item El objeto a representar es muy similar a una superficie.
    \item Solo interesa el aspecto externo del objeto.
    \item El tipo de biblioteca gráfica solo soporta este modelo (no hay más remedio).
\end{itemize}

Para modelar superficies debemos asegurarnos de que las superficies no se intersecan consigo mismas.

Cuando la superficie es cerrada se denominan a los objetos modelos de frontera $\rightarrow$ posibilidad de cambio con modelos sólidos.

Aproximaciones al modelado de superficies:
\begin{itemize}
    \item Discreta y aproximada $\rightarrow$ modelo poliédrico
    \item Continua y exacta $\rightarrow$ modelo matemático
\end{itemize}

\subsubsection{Modelos de Superficie: Poligonales}
Describir una superficie a partir de un conjunto de polígonos conectados mediante vértices y aristas.

Los polígonos se intersecan únicamente en las artistas
\begin{itemize}
    \item Una \textbf{arista} es compartida únicamente por dos polígonos.
    \item Un \textbf{vértice} pertenece a dos o más aristas
\end{itemize}

Cualquier forma 2D o superficie 3D puede aproximarse por polígonos, pero la calidad de la representación mejora a medida que aumentamos el número de polígonos.

\textbf{Ventajas}
\begin{itemize}
    \item Simplicidad 
    \item Fácil y rápido de renderizar 
    \item Ocupa poca memoria
\end{itemize}

\textbf{Limitaciones}
\begin{itemize}
    \item La naturaleza no es poligonal.
    \item Aliasing de polígonos, alisamiento para quitar bordes de sierra, debido a la discretización.
    \item No es sencillo modelar objetos complejos.
\end{itemize}

Normalmente, se trabaja a partir de primitivas poligonales optimizadas, por su reducción de la trasferencia de datos, el problema es la perdida de flexibilidad.
\begin{itemize}
    \item Tira de cuadrados / Tira de triángulos
    \item Matriz de cuadrado / Abanico de triángulos
\end{itemize}

La primitiva poligonal más habitual es el \textbf{Triangulo}, es siempre convexo, matemáticamente simple y siempre coplanares. Además, cualquier polígono puede ser descompuesto en triángulos.

PovRay utiliza triángulos para definir superficies: Mesh almacena eficientemente un número de triángulos y Mesh2 utilizado para la conversión entre formatos gráficos. 

\subsubsection{Modelos de Superficie: Matemáticos}
Ecuaciones matemáticas describen las superficies
\begin{itemize}
    \item \textbf{Superficies Implícitas} $F( x, y, z) = 0$. 
    \begin{itemize}
        \item \textbf{Cuádricas}
        
        Primitivas matemáticas que responden a la ecuación, que dependiendo de los coeficientes será una forma u otra:
        $$a\cdot x^2+b\cdot y^2+c\cdot z^2+2d\cdot xy+2e\cdot yz+2f\cdot xz+2g\cdot z+2h\cdot y+2j\cdot z+k=0$$
        Limitada variedad de las formas (no es simple extenderlo a representación de sólidos) y no sirven para modelar elementos naturales.

        \item \textbf{Superficies equipotenciales o Isosuperficies}
        
        Útiles para representación de formas suaves

        Cada superficie queda definida por el conjunto de puntos con un determinado campo, se modela según el equilibrio entre campos de potencia (equipotencial).

        Al interactuar objetos se crean superficies suaves. El modelado se basa en mover los puntos o elementos que generan el campo, habitualmente los campos son esféricos o cilíndricos.

        \item \textbf{Blobs:} Globulares de formas abultadas.
        $$f(x,y,z)=\Sigma_k b_k e^{-r_k^2a_k}-T=0$$
        T = punto iniciar. a = abultamiento (positivo). b = hendidura (negativos).
        \item \textbf{Metaballs y Softobjects:} Bultos gaussianos
        $$f(r)=\begin{cases}b(1-3r^2/d^2) & si\; r \in [0,d/3]\\ \frac 3 2 b(1-r/d)^2 & si\; r \in [d/3,d] \\ 0 & si\; r>d\end{cases}$$
    \end{itemize} 
    \item \textbf{Superficies Paramétricas} $( x, y, z) = f( u, v)$. Bezier, B-splines y NURBS.
    
    Características
    \begin{itemize}
        \item Facilidad de modelado de formas libres: Control local de deformaciones y Posibilidades de cálculo de tangencias y curvatura
        \item Buenas prestaciones en transformaciones a modelos poligonales 
        \item Buenas prestaciones en cuanto a almacenamiento
    \end{itemize}
    Dos aproximaciones
    \begin{enumerate}
        \item Con superficie que pasa por un conjunto de puntos $\rightarrow$ \textbf{Superficies de Interpolación}
        \item Un conjunto de puntos controlan la forma de la superficie, que no pasa necesariamente por ellos $\rightarrow$ \textbf{Superficies de aproximación}
    \end{enumerate}
    Podemos aproximar cualquier curva unidimensional
    \begin{itemize}
        \item La más básica se basa en el uso de funciones polinómicas $C(u)=\Sigma^n_{i=0} a_iu^i$
        \item La curva depende del grado del polinomio: para definir una curva con n puntos es necesario un polinomio de grado n (problema!!!)
        \item Otra posibilidad consiste en definir la curva como combinación de trozos de curvas: splines. Cuidando los grados de continuidad.
    \end{itemize}
    Tipos:
    \begin{itemize}
        \item \textbf{Splines naturales} 
        
        Normalmente se basan en polinomios de grado cúbico (n=3 en la ecuación), donde se exige continuidad y diferenciabilidad doble.

        Por tanto, la curva consiste en n-1 trozos de grado 3, donde n es el número de puntos por los que pasa la curva.
        \item \textbf{Curvas de Bézier}
        
        Solamente pasan por los puntos extremos y cumplen la propiedad de cierre convexo del polígono de control.
        
        Se definen a partir de la combinación de una familia especial de polinomios: los polinomios de Berstein

        Dos aproximaciones: Curvas de Bezier basadas en polinomios de grado n y de las definidas a trozos.

        Curva de Bezier generales: \\$C(u)=\Sigma^n_{i=0}B_{i,n}(u)P_i \;\; 0 \leq u \leq 1$ con \\$B_{i,n}(u)=\frac{n!}{i!(n-1)!} u^i(1-u)^{n-i}$  con $\Sigma^n_{i,n}(u)=1 \forall u$
        \item \textbf{B-splines}
        
        Similar a las splines naturales, pero sin interpolación. Se une el espacio paramétrico de los diferentes trozos. Esto se realiza definiendo un vector de nudos.
        
        Ventajas
        \begin{enumerate}
            \item Se puede seleccionar el grado de los polinomios base para controlar la suavidad de la curva (grado d)

            \item Permiten control local sobre la forma de la superficie
        \end{enumerate}
        $C(u)=\Sigma^n_{i=0}N_{i,d}(u)P_i$ $u_{min}<u<u_{max}, 2 \leq d \leq n+1$

        $N_{i,1}(u)=$ 1 si $u_i < u < u_{i+1}$ en caso contrario $0$.

        $N_{i,d}(u)= \frac{u-u_i}{u_{i+d-1}-u_i}N_{i,d-1}(u)+\frac{u_{i+d}-u}{u_{i+d}-u_{u+1}}N_{i+1,d-1}(u)$

        Según la distribución de vectores de nudos:
        \begin{itemize}
            \item \textbf{B-splines Racionales} Se definen a partir de la razón o promedio de los polinomios base. Se pueden representar cuádricas y otras superficies de manera exacta.
            $$C(u)=\frac{\Sigma_{i=0}^n w_i N_{i,d}(u)P_i}{\Sigma_{i=0}^n w_i N_{i,d}(u)}$$
            \item \textbf{NURBS-Non Uniform Rational B-Splines:} Cuando podemos tener cualquier distribución del vector de nudos.
        \end{itemize}     
    \end{itemize}
    \item \textbf{Superficies Explícitas} $Z = f( x, y)$, se tratan como superficies paramétricas
\end{itemize}

\subsection{Modelado Sólido}
\subsubsection{Modelos de Barrido}
\begin{enumerate}
    \item Definición de una figura bidimensional (por sus bordes)
    \item Desplazamiento de la figura bidimensional a lo largo de un camino y los bordes generan la superficie.
    \begin{itemize}
        \item Rotación respecto de un eje “Superficie de revolución” 
        \item Desplazamiento a lo largo de una línea recta “Extrusión”
    \end{itemize}
\end{enumerate}
\subsubsection{Geometría Sólida Constructiva (GSC)}
La combinación de primitivas u objetos mediante operaciones booleanas, produciendo objetos más complejos. 

Los objetos diseñados con CSG se representa mediante un árbol binario con nodos y operaciones 

Las operaciones básicas son: Unión, Diferencia, Intersección y Fusión.
\chapter{Transformaciones Geométricas}
\section{Introducción}
\textbf{Transformación geométrica:} Variación del tamaño, forma, posición y orientación de un objeto dentro de una escena.
\begin{itemize}
    \item Transformaciones geométricas afines en 2D y 3D.
    \begin{itemize}
        \item Transformaciones lineales y traslaciones:
        \begin{itemize}
            \item Traslación
            \item Escalamiento 
            \item Rotación
        \end{itemize}
    \end{itemize}
    \item Transformaciones no afines.
\end{itemize}

\section{Aplicaciones en Informática Gráfica}
Para modelar objetos.
\begin{itemize}
    \item Ejem. Blender se puede partir de un cubo, escalarlo en una dirección y rotarlo para formar una columna.
\end{itemize}

Para crear objetos complejos.

Para animar objetos, traslaciones y rotaciones son los elementos básicos del movimiento.

\section{Transformaciones lineales}
Una transformación es una función: $v(x,y,z) \rightarrow v'(x',y',z')$.

Una transformación es lineal, si se cumple:
\begin{itemize}
    \item $T(u+v) = T(u) + T(v)$ La transformación del vector suma es la suma de las transformaciones.
    \item $T(x\cdot u) = x\cdot T(u)$, siendo x un escalar.
\end{itemize}

\textbf{Son lineales:} Rotación, Escalado y Estiramiento (shear).

En las transformaciones lineales se cumple que las líneas paralelas siguen siendo paralelas (los ángulos no tienen por qué mantenerse).

El origen de coordenadas se mantiene.

\textbf{No es lineal:} Traslación.

En la traslación se mantiene el paralelismo, pero el origen de coordenadas no se mantiene.

La traslación no es una transformación lineal, pero sí es una transformación afín.

\textbf{Transformación afín:} transformación lineal y traslación.

\section{Transformaciones 2D}
\subsection{Forma General}
Las transformaciones lineales se pueden representar con matrices cuadradas. La traslación también, pero usando coordenadas homogéneas (una dimensión más).

$$\left[\begin{matrix}
x' \\ y'
\end{matrix}\right] =
\left[\begin{matrix}
a & b \\ c & d
\end{matrix}\right]
\left[\begin{matrix}
x \\ y
\end{matrix}\right]=
\left[\begin{matrix}
ax+by \\ cx+dy
\end{matrix}\right]$$
\begin{itemize}
    \item $a$ es eje x.
    \item $d$ es eje y.
    \item $b$ es deformación en x.
    \item $c$ es deformación en y.
\end{itemize}

\subsection{Transformación lineal 2D}
Vector como combinación lineal de los vectores base $i$, $j$. Cuando se haga una transformación cambia el vector base.
\begin{itemize}
    \item Ejem. $v (-1,2) = -1 i + 2 j$
\end{itemize} 
$i$ es la primera columna de la matriz de transformación y $j$ la segunda.

\subsection{Traslación}
Cambio de la posición de un objeto a lo largo de una línea recta.

Desplaza cada punto de una figura o espacio la misma cantidad en una determinada dirección, vector de traslación $(d_x,d_y)$. $d_x$ unidades paralelas al eje x y $d_y$ unidades paralelas al eje y.

$$\left[\begin{matrix}
x' \\ y'
\end{matrix}\right] =
\left[\begin{matrix}
x \\ y
\end{matrix}\right]+
\left[\begin{matrix}
d_x \\ d_y
\end{matrix}\right]=
\left[\begin{matrix}
x+d_x \\ y+d_y
\end{matrix}\right]$$

NO es una transformación lineal, aunque si afín. La transformación de la suma de dos vectores no es la suma de las transformaciones.

\subsection{Escalado}
Alteración de las dimensiones de un objeto:
\begin{itemize}
    \item Segmentos paralelos al eje x se multiplican por $s_x$.
    \item Segmentos paralelos al eje y se multiplican por $s_y$.
\end{itemize}

Cuando $s_x = s_y$ es un escalado uniforme.

$$\left[\begin{matrix}
x' \\ y'
\end{matrix}\right] =
\left[\begin{matrix}
s_x & 0 \\ 0 & s_y
\end{matrix}\right] \cdot
\left[\begin{matrix}
x \\ y
\end{matrix}\right]=
\left[\begin{matrix}
x\cdot s_x \\ y\cdot s_y
\end{matrix}\right]$$

\subsection{Rotación}
Cambiar la orientación de un objeto $\theta$ respecto del origen.

$\phi$ es la posición angular del punto $P(x, y)$ respecto al origen.

Los vectores base pasarán a ser:
\begin{itemize}
    \item $i'=(1 \cos \theta , 1 \sin \theta)$
    \item $j'=(-1 \sin \theta, 1 \cos \theta)$
\end{itemize}

$$\left[\begin{matrix}
x' \\ y'
\end{matrix}\right] =
\left[\begin{matrix}
\cos \theta & - \sin \theta \\ \sin \theta & \cos \theta
\end{matrix}\right] \cdot
\left[\begin{matrix}
x \\ y
\end{matrix}\right]=
\left[\begin{matrix}
x \cos \theta - y \sin \theta \\ x \sin \theta + y \cos \theta
\end{matrix}\right]$$

\subsection{Cuidado con el orden de las transformaciones}
Secuencia Rotación y Escalado no es equivalente a Escalado y Rotación.

Se multiplica de Derecha a Izquierda.

No se cumple la propiedad conmutativa, pero si la asociativa.
\begin{itemize}
    \item $A2 = S \cdot R \cdot A$
    \item $A2 = (S \cdot R) \cdot A$
\end{itemize}

Aunque existen casos especiales donde al aplicar transformaciones de traslación, escalado y rotación, el producto de matrices es conmutativo, como:
\begin{itemize}
    \item Traslación seguida de traslación.
    \item Escalado seguido de escalado
    \item Rotación seguida de rotación
    \item Rotación y escalado uniforme
\end{itemize}

Se busca ejecutar la secuencia de transformaciones más eficientemente, para esto necesitamos un sistema de referencia que sea homogéneo y entre la traslación.

\section{Sistema de Coordenadas Homogéneas 2D}
Un punto $( x , y )$ en el sistema cartesiano se representa por $( xW , yW , W )$ en el sistema homogéneo, donde $W\neq 0$ puede ser cualquier número real (geometría proyectiva).

Cuando $W\neq 1$, para simplificar cálculos, dividimos todas las coordenadas por W, obteniendo $( x , y , 1 )$ (W suele utilizarse en el eje z profundidad).

Cuando $W = 0$, estamos realizando transformaciones del punto $( x , y )$ sobre el plano $z = 1$.

Las matrices que representan puntos en 2D serán de 3x3.

\subsection{Traslación en coordenadas homogéneas}
Ahora la traslación podrá hacerse representarse como una matriz.

$$\left[\begin{matrix}
x' \\ y' \\ 1
\end{matrix}\right] =
\left[\begin{matrix}
1 & 0 & d_x \\ 0 & 1 & d_y \\ 0 & 0 & 1
\end{matrix}\right] \cdot
\left[\begin{matrix}
x \\ y \\ 1
\end{matrix}\right]=
\left[\begin{matrix}
x + d_x \\ y + d_y \\ 1
\end{matrix}\right]$$

Es como el bias de las redes de neuronas, el 1 se multiplica por un desplazamiento.

Ahora se puede hacer composición de traslaciones.

$P'=T(d_{x1},d_{y1}) \cdot P; \;\;P''=T(d_{x2},d_{y2}) \cdot P'$

$P''=T(d_{x2},d_{y2}) \cdot (T(d_{x1},d_{y1}) \cdot P)=(T(d_{x2},d_{y2}) \cdot T(d_{x1},d_{y1})) \cdot P$

$$\left[\begin{matrix}
1 & 0 & d_{x1} \\ 0 & 1 & d_{y1} \\ 0 & 0 & 1
\end{matrix}\right] \cdot
\left[\begin{matrix}
1 & 0 & d_{x2} \\ 0 & 1 & d_{y2} \\ 0 & 0 & 1
\end{matrix}\right]=
\left[\begin{matrix}
1 & 0 & d_{x1}+d_{x2} \\ 0 & 1 & d_{y1}+d_{y2} \\ 0 & 0 & 1
\end{matrix}\right]$$

\subsection{Escalado en coordenadas homogéneas}
$$\left[\begin{matrix}
x' \\ y' \\ 1
\end{matrix}\right] =
\left[\begin{matrix}
s_x & 0 & 0 \\ 0 & s_y & 0 \\ 0 & 0 & 1
\end{matrix}\right] \cdot
\left[\begin{matrix}
x \\ y \\ 1
\end{matrix}\right]=
\left[\begin{matrix}
x\cdot s_x \\ y\cdot s_y \\ 1
\end{matrix}\right]$$

$$\left[\begin{matrix}
    s_{x2} & 0 & 0 \\ 0 & s_{y2} & 0 \\ 0 & 0 & 1
\end{matrix}\right] \cdot
\left[\begin{matrix}
    s_{x1} & 0 & 0 \\ 0 & s_{y1} & 0 \\ 0 & 0 & 1
\end{matrix}\right]=
\left[\begin{matrix}
s_{x1} \cdot s_{x2} & 0 & 0 \\ 0 & s_{y1} \cdot s_{y2} & 0 \\ 0 & 0 & 1
\end{matrix}\right]$$

\subsection{Rotación en coordenadas homogéneas}
$$\left[\begin{matrix}
x' \\ y' \\ 1
\end{matrix}\right] =
\left[\begin{matrix}
\cos \theta & - \sin \theta & 0 \\ \sin \theta & \cos \theta & 0 \\ 0 & 0 & 1
\end{matrix}\right] \cdot
\left[\begin{matrix}
x \\ y \\ 1
\end{matrix}\right]=
\left[\begin{matrix}
x \cos \theta - y \sin \theta \\ x \sin \theta + y \cos \theta \\ 1
\end{matrix}\right]$$

$$\left[\begin{matrix}
\cos \alpha & - \sin \alpha & 0 \\ \sin \alpha & \cos \alpha & 0 \\ 0 & 0 & 1
\end{matrix}\right] \cdot
\left[\begin{matrix}
\cos \beta & - \sin \beta & 0 \\ \sin \beta & \cos \beta & 0 \\ 0 & 0 & 1
\end{matrix}\right]=
\left[\begin{matrix}
\cos (\alpha + \beta ) & - \sin (\alpha + \beta ) & 0 \\ \sin (\alpha + \beta ) & \cos (\alpha + \beta ) & 0 \\ 0 & 0 & 1
\end{matrix}\right]$$

\subsection{Transformaciones inversas}
\textbf{Traslación} $T( d_x, d_y)$ es: $T^{-1} ( d_x , d_y ) = T( -d_x , -d_y )$

\textbf{Escalado} $S( s_x , s_y )$ es: $S^{-1} ( s_x , s_y ) = S(1/s_x ,1/s_y )$

\textbf{Rotación} $R(T)$ es: $R^{-1}(\theta) = R(-\theta)$ que equivale a cambiar los signos de los senos de la transformación normal.

\subsection{Transformaciones generales afines}
Producto de una secuencia arbitraria de matrices de rotación, traslación y escalamiento.

Propiedad de conservar el paralelismo de las líneas, pero no longitudes ni ángulos. Rotaciones, escalamientos y traslaciones subsiguientes no podrían hacer que las líneas dejen de ser paralelas (nótese que sí permiten deformar el objeto).

\subsection{Estiramiento (Shear)}
El estiramiento puede realizarse respecto de cualquier eje.

$$SH_x = \left[\begin{matrix}
1 & a & 0 \\ 0 & 1 & 0 \\ 0 & 0 & 1
\end{matrix}\right] \;\;\;
SH_y = \left[\begin{matrix}
1 & 0 & 0 \\ b & 1 & 0 \\ 0 & 0 & 1
\end{matrix}\right]$$

\subsection{Composición de transformaciones}
El objetivo es ganar eficiencia aplicando una sola transformación compuesta a un punto, como es la aplicación de una rotación a un objeto sobre un punto P, que no es el origen.
\begin{enumerate}
    \item Traslación al origen $T(-x_1, -y_1)$
    \item Rotación $R(\theta)$
    \item Traslación al punto P $T(x_1, y_1)$
\end{enumerate}

\subsection{Operaciones respecto a un punto $P_1(x_1,y_1)$}
Recordar que se hace las operaciones de derecha a izquierda.

\subsubsection{Rotación}
$$T(x_1, y_1)\cdot R(\theta) \cdot T(x_1, y_1) = 
\left[\begin{matrix}
\cos \theta & - \sin \theta & x_1(1-\cos \theta)+y_1 \sin \theta \\ \sin \theta & \cos \theta & y_1(1-\cos \theta)-x_1 \sin \theta \\ 0 & 0 & 1
\end{matrix}\right]$$

\subsubsection{Escalado}
$$T(x_1, y_1)\cdot S(s_x, s_y) \cdot T(x_1, y_1) = 
\left[\begin{matrix}
s_x & 0 & x_1(1-s_x) \\ 0 & s_y & y_1(1-s_y) \\ 0 & 0 & 1
\end{matrix}\right]$$

\subsubsection{Escalado y Rotación}
$M=T(x_1, y_1) \cdot R(\theta) \cdot S(s_x, x_y) \cdot T(-x_1, -y_1)$

\section{Sistema de Coordenadas Homogéneas 3D}
Un punto $(x, y, z)$ en el sistema cartesiano se representa por $( xW , yW , zW, W )$ en el sistema homogéneo, donde $W\neq 0$ puede ser cualquier número real.

Cuando $W\neq 1$, para simplificar cálculos, dividimos todas las coordenadas por W, obteniendo $( x , y , z, 1 )$

\subsection{Traslación}
$$T(d_x, d_y, d_z) \cdot
\left[\begin{matrix}
x \\ y \\ z \\ 1
\end{matrix}\right] = 
\left[\begin{matrix}
1 & 0 & 0 & d_x \\ 0 & 1 & 0 & d_y \\ 0 & 0 & 1 & d_z \\ 0 & 0 & 0 & 1
\end{matrix}\right] \cdot
\left[\begin{matrix}
x \\ y \\ z \\ 1
\end{matrix}\right] = 
\left[\begin{matrix}
x+ d_x \\ y+d_y \\ z+d_z \\ 1
\end{matrix}\right]$$

\subsection{Escalado}
$$S(s_x, s_y, s_z) \cdot
\left[\begin{matrix}
x \\ y \\ z \\ 1
\end{matrix}\right] = 
\left[\begin{matrix}
s_x & 0 & 0 & 0 \\ 0 & s_y & 0 & 0 \\ 0 & 0 & s_z & 0 \\ 0 & 0 & 0 & 1
\end{matrix}\right] \cdot
\left[\begin{matrix}
x \\ y \\ z \\ 1
\end{matrix}\right] = 
\left[\begin{matrix}
x\cdot s_x \\ y\cdot s_y \\ z\cdot s_z \\ 1
\end{matrix}\right]$$

\subsection{Rotación}
$$R_x(\theta)=\left[\begin{matrix}
1 & 0 & 0 & 0 \\ 
0 & \cos \theta & - \sin \theta & 0 \\ 
0 & \sin \theta & \cos \theta & 0 \\ 
0 & 0 & 0 & 1
\end{matrix}\right]
R_y(\theta)=\left[\begin{matrix}
\cos \theta & 0 & \sin \theta & 0 \\ 
0 & 1 & 0 & 0 \\ 
- \sin \theta & 0 & \cos \theta & 0 \\ 
0 & 0 & 0 & 1
\end{matrix}\right]
R_z(\theta)=\left[\begin{matrix}
\cos \theta & -\sin \theta & 0 & 0 \\ 
\sin \theta & \cos \theta & 0 & 0 \\ 
0 & 0 & 1 & 0 \\
0 & 0 & 0 & 1
\end{matrix}\right]$$

\subsection{Rotación en coordenadas homogéneas}
La composición de una secuencia arbitraria de rotaciones con respecto a los ejes x, y, z como submatriz 3x3 ortogonal. La inversa de una matriz ortogonal es su transpuesta.
$$R_z(\theta)=\left[\begin{matrix}
r11 & r12 & r13 & 0 \\ 
r21 & r22 & r23 & 0 \\ 
r31 & r32 & r33 & 0 \\
0 & 0 & 0 & 1
\end{matrix}\right]$$
\pagebreak

\section{Rotación alrededor de cualquier eje}
\begin{itemize}
    \item Primero movemos el eje hasta el eje z.
    \begin{enumerate}
        \item Rotamos alrededor de x $R_x$
        \item Rotamos alrededor de y $R_y$
    \end{enumerate}
    \item Después, aplicamos la rotación deseada seguida de las inversas usadas para colocarlo sobre el eje z.
    \begin{enumerate}
        \item Rotamos lo deseado $R_z$
        \item Deshacemos la rotación $R_x^{-1}R_y^{-1}$
    \end{enumerate}
\end{itemize}
$$R_{final}=R_x^{-1}R_y^{-1}R_zR_yR_x$$

\section{Rotación alrededor de cualquier eje y punto de origen}
\begin{enumerate}
    \item Traslación al origen deseado $T_{xyz}$
    \item Rotamos $R_x^{-1}R_y^{-1}R_zR_yR_x$
    \item Invertir traslación $T_{xyz}^{-1}$
\end{enumerate}
$$R_{final}=T_{xyz}^{-1}R_x^{-1}R_y^{-1}R_zR_yR_xT_{xyz}$$

\section{Matrices en PovRay}
Las matrices se definen con matrix <v00, v01, v02, v10, ...>, en el que la última columna se asume que es 0, 0, 0, 1. De manera que la matriz está traspuesta, siendo la última fila la de traslación 0, 0, dy, dx (está rotada).
\chapter{Primitivas 2D}
\section{Definiciones}
\textbf{Pixel:} Mínima unidad homogénea de color (intensidad) de una imagen digital. Área de la pantalla que tiene asociada una posición de memoria (pixel de 8 bits tiene 256 variaciones de color).

\textbf{Primitivas de salida:} Estructuras geométricas básicas continuas (puntos, rectas, curvas, áreas, ...) empleadas para generar imágenes

\textbf{Conversión al ''raster'':} Aproximar primitivas mediante un conjunto de puntos discreto (mapa de píxeles)
\begin{itemize}
    \item Localizar posiciones de pixel más próximas al objeto.
    \item Almacenados valores intensidad/color en búfer (PIXMAP)
    \item Trazado de imagen en pantallas (intensidad a píxeles)
\end{itemize}

\textbf{Búfer de imagen:} Área de memoria que almacena el conjunto de valores de intensidad para todos los puntos de la pantalla.
\begin{itemize}
    \item B/N 1 bit por pixel, 2 intensidades. BITMAP.
    \item Color n bits por pixel, $2^n$ intensidades. PIXMAP.
    \item Sistemas alta calidad hasta 24 bits por pixel. Sistema color RGB de 24 bits.
\end{itemize}

\textbf{Procesador graficó/ Controlador de gráficas:}  libera CPU de trabajos gráficos.

\section{Conversión de puntos}
Problema - convertir al raster un punto definido mediante una coordenada almacenando el color con el que se va a visualizar.

Algoritmo sencillo de conversión, sea $(x, y)$ el punto real la posición en el raster es $(x', y')=(round(x), round(y))$. $round()$ entero más próximo.

\section{Conversión de rectas}
Problema - calcular las coordenadas de los píxeles que representan una recta infinitamente delgada sobre una malla de un raster 2D.

Un buen rasterizador debe cumplir lo siguiente:
\begin{itemize}
    \item Secuencia de píxeles lo más recta posible
    \item Líneas con el mismo grosor independientemente de la pendiente
    \item El algoritmo debe ser muy rápido
\end{itemize}
 
\subsection{Fuerza bruta}
Ecuación de la recta entre puntos extremos $(x_1 ,y_1)$ y $(x_2 ,y_2 )$, $y=mx+b$.
\begin{itemize}
    \item $m= \frac {(y_2-y_1)}{(x_2-x_1)}$
    \item $b$ se halla usando la pendiente, $m$, y los puntos conocidos. $b=y_1-m*x_1$
\end{itemize}
Para cada punto $x$ entre $x_1$ y $x_2$ se calcula el pixel más cercano a la curva definida por la recta. Se va aumentando progresivamente la $x$.

\textbf{Restricción:} m tiene que estar entre -1 y 1, si el grado que forma es mayor que 45 se forma de manera discontinua.

\textbf{Solución:} Intercambiar la x por la y en los extremos de la regla, hacer el proceso e intercambiarlas de nuevo.

\textbf{Desventajas:} Ineficiente, requiere multiplicaciones en coma flotante y función round().

\subsection{Algoritmo Analizador Diferencial Digital (DDA, Digital Diferenctial Analizer)}

Algoritmo que traza valores sucesivos de $(x, y)$ incrementando simultáneamente $x$ e $y$ con pequeños pasos.

\begin{enumerate}
    \item Calculamos $\Delta x = x_2 -x_1$ y $\Delta y = y_2 -y_1$.
    \item Calculamos los incrementos que usaremos para dibujar, $x_{inc}$ y $y_{inc}$. $x_{inc} = \frac{\Delta x}{\max(|\Delta x|, |\Delta y|)}$ y $y_{inc} = \frac{\Delta y}{\max(|\Delta x|, |\Delta y|)}$.
    \item Para el eje vamos iterando: $x=x+x_{inc}$ y $y=y+y_{inc}$. Redondeando aquellos que sean decimales.
\end{enumerate}

\textbf{Ventaja:} elimina el producto de coma flotante.

\textbf{Desventaja:} invoca la función round.

Seleccionamos un eje con pasos unitarios (según el valor de $|m|$)

Representamos la línea de izquierda a derecha si es positivo y derecha a izquierda  si es negativos.

\subsection{Algoritmo del punto medio de Bresenham (Bresenham midpoint algorithm)}
Consiste en determinar si cuando se calcula el punto debe ir a superior o inferior según la distancia del punto real con estos. El factor de decisión, p, nos indicará cuál coger.

\begin{enumerate}
    \item Calculamos $\Delta x = x_2 -x_1$ y $\Delta y = y_2 -y_1$. Además, para facilitar los cálculos podemos también calcular $2\Delta y - 2 \Delta x$ y $2 \Delta x$.
    \item Calculamos el $p_0=2\Delta y - \Delta x$
    \item Comenzando en $k=0$, repetimos $\Delta x$ veces:
    \begin{enumerate}
        \item Si $p_k < 0$ entonces $x_{k+1}=x_k+1$, $y_{k+1}=y_k$ y $p_{k+1}=p_k+2\Delta y$
        \item Si $p_k > 0$ entonces $x_{k+1}=x_k+1$, $y_{k+1}=y_k+1$ y $p_{k+1}=p_k+2\Delta y - 2 \Delta x$
    \end{enumerate}
    \item Dibujamos los puntos calculados
\end{enumerate}
\chapter{Fractales}
Figuras que se autorrepiten, es decir, la figura está formada de sí misma.

Esta estructura se da en la naturaleza, en los vasos sanguíneos, cuenca del Amazonas o las vías romanas. Se pueden dar en forma de ramificación, en espiral o en conos.
\begin{itemize}
    \item Ejem. Cálculo de la longitud de la costa de Gran Bretaña, que se trata de estimar con un segmento de x metros, cuanto más pequeño lo vayamos haciendo más grande sale la longitud. Por lo que siempre hay un segmento más pequeño. 

\end{itemize}

\section{Definición}
El concepto de fractal aparece en 1975 con el libro Benoit Mandelbrot - ''Los objetos Fractales: forma, azar y dimensión''.

Aunque se habían descubierto antes y se llamaron ''monstruos matemáticos''

Un objeto fractal tiene dos características básicas:
\begin{itemize}
    \item \textbf{Infinito detalle} en cada punto 
    \item \textbf{Auto similitud entre las partes} del objeto y su totalidad
\end{itemize}
Veremos \textbf{siempre la misma figura}, independientemente de lo que nos acerquemos al mismo (fractal auténtico).

Su geometría requiere \textbf{gráficos por ordenador} para visualizarse e investigarlos.

\section{Caracterización}
\begin{table}[H]
    \begin{tabular}{|c|c|}
    \hline
    \rowcolor[HTML]{BFBFBF} 
    \begin{tabular}[c]{@{}c@{}}Geometría\\ Euclídea\end{tabular}                                   & \begin{tabular}[c]{@{}c@{}}Geometría\\ Fractal\end{tabular}                         \\ \hline
    Ecuaciones                                                                                     & \begin{tabular}[c]{@{}c@{}}Procedimientos\\ (sin ecuaciones)\end{tabular}           \\ \hline
    Objetos fabricados                                                                             & Objetos naturales                                                                   \\ \hline
    \begin{tabular}[c]{@{}c@{}}Diferenciable,\\ localmente suave\end{tabular}                      & \begin{tabular}[c]{@{}c@{}}No diferenciable,\\ localmente rugoso\end{tabular}       \\ \hline
    \begin{tabular}[c]{@{}c@{}}Aumenta el detalle\\ y definición a \\ distancia corta\end{tabular} & \begin{tabular}[c]{@{}c@{}}Mismo detalle desde\\  todas las distancias\end{tabular} \\ \hline
    \end{tabular}
    \caption{Comparación Geometría Fractal y Euclídea}
\end{table}

\section{Aplicaciones}
\begin{itemize}
    \item Antenas fractales (ancho de banda)
    \item Compresión de imágenes
    \item Sistemas dinámicos (cotización de valores bursátiles)
    \item Modelado de formas naturales (helechos, montañas o relieve)
    \item Arte fractal.
\end{itemize}

\section{Procedimiento de Generación}
Aplicar recursivamente una función de transformación a los puntos de una región del espacio:
\begin{itemize}
    \item Tenemos una función de transformación $F(X)$ y partimos de punto inicial seleccionado $P_0 =(x_0 , y_0 , z_0 )$
    \item Cada vez que aplicamos F, generamos niveles sucesivos de detalle, como $P_1 =F(P_0 ), P_2 =F(P_1 ), ... P_{k+1} =F(P_k ), ...$
    \item Sistema dinámico realimentado
\end{itemize} 

$F(X)$ puede aplicarse a: Conjunto específico de puntos, o a un Conjunto inicial de primitivas (líneas rectas, áreas de color, superficies o sólidos). 

$F(X)$ puede definirse en términos de transformaciones: Geométricas (escalado, traslación, rotación) o Coordenadas no lineales y parámetros de decisión.

No podemos desplegar variaciones de detalle inferiores a un pixel $\rightarrow$ $F(X)$ se aplica un número finito de veces. Para ver más detalle acercamos la imagen y aplicamos $F(X)$ de nuevo.

Los procedimientos a utilizar pueden ser Deterministas o Aleatorios / estocásticos.

\section{Clasificación de Fractales}
\begin{itemize}
    \item \textbf{Autosimilares:} Sus partes son versiones a escala de su totalidad. Se aplica un parámetro de escala (S) toda la forma inicial.
    
    \textbf{Autosimilar estadísticamente} si aplicamos variaciones aleatorias.
    \item \textbf{Autoafines:} Un factor de escala en cada eje, $S_x, S_y, S_z$.
    
    \textbf{Autoafines estadísticamente} si aplicamos variaciones aleatorias.
    \item \textbf{Invariantes:} Transformaciones no lineales. Fractales autocuadráticos (Mandelbrot), autoinversos...
\end{itemize}

\section{Dimensión Fractal}
Medida de la variación de detalle de un objeto fractal.

En la geometría euclídea los objetos tienen dimensión entera. Los objetos en la geometría fractal tiene dimensión fraccionaria (o dimensión fractal).

La dimensión fractal de un objeto siempre es mayor que su dimensión euclídea (o topológica).

Dada una dimensión, podemos generar objetos fractales mediante procedimientos recursivos.

Podríamos calcular la dimensión fractal de un objeto a partir de las propiedades del mismo, mediante otros procedimientos (tarea difícil de realizar).

Cálculo de la dimensión de un fractal autosimilar determinista (un solo factor de escala) por analogía a las subpartes de un objeto euclidiano.

$NS^{DE}=1 \rightarrow NS^{D}=1$
\begin{itemize}
    \item DE = Dimensión euclídea
    \item D = Dimensión fractal
    \item N = n. piezas semejantes
    \item S = factor de escala
\end{itemize}
$$N=\frac{1}{S^D}=\left(\frac{1}{S}\right)^D; \ln N = D \ln \frac{1}{S}; D = \frac{\ln N}{\ln \frac{1}{S}}$$
\chapter{Teoría del Color}
\section{Introducción}
La percepción del color de los objetos depende de:
\begin{itemize}
    \item Propiedades del objeto
    \item Fuente de luz
    \item Color del entorno
    \item Sistema visual humano (no solo el ojo, también el cerebro)
\end{itemize}

La realidad física, son longitudes de onda, pero la percepción son los colores.

Existen diversas teorías, técnicas de medida, estándares..., pero no hay una teoría universalmente aceptada sobre la percepción del color que realizamos los humanos.

La ciencia del color es una mezcla de componentes subjetivos (percepción) con componentes físicos (longitudes de onda, espectros)

\section{La Luz}
La luz (gama de colores) es la banda de longitudes (o frecuencias) de onda del espectro electromagnético, que es visible para nosotros. Son de longitud 750-400 nm.

\subsection{Historia}
Siglo XIX aceptaban que la luz tenía naturaleza ondulatoria, aunque no se relacionaba con la electricidad ni con el magnetismo.
\begin{itemize}
    \item Pero todas las ondas conocidas necesitaban un medio para propagarse y la luz viajaba en el vacío con más velocidad que en medios como el aire o el agua. Se proponía la existencia del éter luminífero, que se desecha en el s. XIX.
\end{itemize}

En 1845 Faraday (aproximación práctica) descubrió que el magnetismo altera el plano de polarización de la luz (relación entre la luz y el magnetismo). Y especuló que la luz podía ser resultado de la vibración de las líneas de fuerza eléctricas y magnéticas. 
\pagebreak

Nadie le hizo caso, hasta que Maxwell (aproximación teórica) publica las ecuaciones del electromagnetismo. Deduce que deben existir perturbaciones electromagnéticas que se autopropagan en el vacío. Mediante una ecuación consigue sin necesidad de punto de referencia aproximar la velocidad de la luz mediante campos.

Quedó aceptado que la luz es una radiación electromagnética.

A finales del siglo XIX se encontraron nuevos efectos que no se podían explicar con la naturaleza ondulatoria de la luz.
\begin{itemize}
    \item Radiación del cuerpo negro
    \item Efecto fotoeléctrico
\end{itemize}

Solo se podían explicar si la luz consistía en partículas: cuantos de energía, fotones. Plank $E=hv$

La luz se puede manifestar como onda o como partícula (dualidad ondapartícula)

\subsection{Radiación de cuerpo negro}
La energía que emite un cuerpo depende de su temperatura. Esta temperatura les confiere un color que puede parecer contradictorio, son:

Temperatura más alta $\rightarrow$ Temperatura de color frío (azulado)

Temperatura más baja $\rightarrow$ Temperatura de color cálido (rojizo)

El sol actúa como una especie de cuerpo negro, cuya radiación solar en el camino se ve afectada, de manera que lo que se ve fuera de la tierra es diferente que en la superficie.

\section{Sistema visual humano}
Tres tipos de conos: S (corto), M (medio),L (largo). Cada uno responde a diferentes rangos de longitudes de onda. Estos se alinean bastante bien con el azul, verde y rojo.

Nuestro cerebro con 3 señales interpreta cualquier color que podamos percibir.

Las aves tienen cuatro tipos de conos, uno de ellos detecta la luz UV

Puede ocurrir que percibamos igual diferentes espectros, mezclas de colores.
\begin{itemize}
    \item La luz espectral es monocromática, y no se puede descomponer en otros.
    \item Luz resultado de mezclar colores, se puede con un prisma descomponer en los colores que lo componen.
\end{itemize}

\section{Caracterización del Color}
Cada color del espectro puede describirse en función de su frecuencia de onda ($f$, medida en hercios) o su longitud de onda ($\lambda$, medida en nanómetros), que están relacionadas por $c = \lambda f$.

\textbf{Luz monocromática o Color espectral:} está formada por una sola longitud de onda (colores del arcoíris). Ejem. El rosa no existe como onda.

\textbf{Fuente de luz blanca (sol)} emite todas las frecuencias. Al incidir sobre un objeto, este absorbe algunas longitudes y refleja otras. 

\textbf{La combinación de longitudes de onda reflejadas determina el color (tono, matiz) del objeto}

Desde el punto de vista de nuestra percepción, un color se caracteriza con 3 parámetros (\textbf{modelo HSB}/HSL/HSV):
\begin{itemize}
    \item \textbf{Hue} (Tono, matiz, color): El color, longitud de onda predominante, se puede discriminar aproximadamente entre 150 tonos.
    \item \textbf{Saturation} (Saturación): Es el grado de pureza del color observado: va desde el color puro al gris. Los menos saturados contienen más luz blanca.
    \item \textbf{Brightness}, Lightness, Value (Brillo): Intensidad de la luz que se percibe. Cuanta energía emite.
    
    Intensidad - energía radiante que se emite por unidad de tiempo
\end{itemize}

Podemos distinguir del orden de 7 millones de colores

\textbf{Colores primarios:} Aquellos que no pueden obtenerse mediante la mezcla de ningún otro.
\begin{itemize}
    \item Colores aditivos - al añadirse generan el blanco. Azul, verde y rojo.
    \item Colores substractivos - filtros que absorben la luz. Cian, amarillo y magenta.
\end{itemize}

\textbf{Colores complementarios:} Combinados producen luz blanca. Rojo-cian, verde-magenta, azul-amarillo

Se pueden combinar fuentes de luz de color con distinta intensidad para, seleccionando adecuadamente las intensidades, generar un determinado rango de colores.

\textbf{No se puede} combinar un conjunto finito o real de colores primarios para \textbf{generar todos los colores visibles} posibles

Modelos de color emplean \textbf{3 colores primarios para obtener la gama de colores (gamut)} asociada a ese modelo.

\section{Estandarización}
Necesidad de estandarizar los colores percibidos. En principio, con tres valores podríamos representar cualquier color.

\textbf{Color matching experiment:} Mezcla de tres colores primarios (ej. azul, verde y rojo del modelo RGB) para generar todos los colores del espectro visual.
\begin{itemize}
    \item Combinar colores puros RGB para obtener un cierto color espectral.
    \item Se hace para todos los colores del espectro y con diferentes personas.
    \item Lo que si varía de los colores base es la intensidad de los colores (ya que el matiz es propio de la longitud de RGB)
    \item No todos los colores del espectro pueden ser igualados mediante la mezcla de R, G y B como el cian de 500 nm, pero si se le puede restar rojo.
    \item \textbf{Resultado:} Para cada color del espectro, obtenemos los valores RGB necesarios.
\end{itemize}

Estos experimentos son la base de la ciencia del color. Dieron lugar al espacio de color \textbf{CIE 1931} establecido por la \textbf{Comission Internationale de l'Eclairage (CIE)}.

\section{Colores Primarios Estándar - Modelo CIE XYZ}
\textbf{La Comisión Internacional sobre Iluminación, CIE (1931)} define un estándar de tres colores imaginarios aditivos que no necesiten valores negativos. 

Se hace una transformación lineal del espacio R,G,B en el espacio X,Y,Z (vectores en espacio 3D de color aditivo) para que todos los valores de estos nuevos colores primarios sean positivos. $\left(\begin{matrix} X \\ Y \\ Z \end{matrix}\right)=M\left(\begin{matrix} R \\ G \\ B \end{matrix}\right)$. Combinando los colores imaginarios XYZ se puede generar todo el espectro.

Se eligen los primarios de forma que Y es igual a la luminosidad.

\section{CIE 1931: diagrama de cromaticidad (2D)}
\textbf{Colores imaginarios XYZ:} no son RGB, XYZ son combinaciones lineales de los valores de RGB para que todos los valores sean positivos.

Para hacerlo más manejable, se proyectará el espacio XYZ (que es 3D) en un plano xy $\rightarrow$ esto es el diagrama de cromaticidad (también llamado CIE xyY)

Esto se consigue normalizando los valores X, Y, Z, para que su suma sea 1 ($x+y+z=1$). $x=\frac{X}{Z+Y+Z}$ $y=\frac{Y}{Z+Y+Z}$ $z=\frac{Z}{Z+Y+Z}$.

El color puede dividirse en dos partes: brillo y cromaticidad.
\begin{itemize}
    \item $Y$ es una medida del brillo o luminosidad de un color. Máxima en el amarillo verdoso
    \item La cromaticidad se determina a través de los dos parámetros derivados $x$ e $y$ (dos
    de los valores de X,Y, Z)
\end{itemize}

\section{Diagrama de Cromaticidad}
Proyección del plano $x+y+z=1$ en el plano XY.
\begin{itemize}
    \item A partir de la cromaticidad $x$ e $y$ (fija relación entre colores) y la luminancia $Y$, se obtienen los valores mediante: $X=x\frac{Y}{y}$ y $Z=z\frac{Y}{y}$
\end{itemize}

El diagrama representa toda la cromaticidad percibida por el ojo humano
\begin{itemize}
    \item Aplicaciones: Nombrar y normalizar colores, Mezcla de colores, Identificar colores complementarios, Determinar la frecuencia dominante y Comparar gamut de dispositivos.
    \item Propiedades
    \begin{itemize}
        \item Los colores del borde son: \textbf{Colores espectrales puros (spectral locus)}
        \item La línea recta es: \textbf{La línea del púrpura (purple boundary)}
        \item La suma de dos colores se encuentra en la línea que los une
        \item El blanco se encuentra en $x=1/3$, $y=1/3$. La línea que une dos colores complementarios pasa por ese punto.
    \end{itemize}
    \item Se puede representar cualquier color con:
    \begin{itemize}
        \item los distintos valores del \textbf{triestímulo XYZ}
        \item o con los valores de \textbf{luminancia Y y de cromaticidad x,y}
    \end{itemize}
\end{itemize}

Mezcla de colores
\begin{itemize}
    \item Mezclando los \textbf{colores I,J} se pueden obtener toda la \textbf{gama} situada en el \textbf{segmento $\overline{IJ}$}
    \item Con los colores \textbf{I,J,K se obtienen los colores del triángulo inscrito (gamut)}
\end{itemize}
\pagebreak

Color blanco
\begin{itemize}
    \item El blanco es una mezcla de todos los colores.
    \item Se define un blanco estándar, blanco de igual energía.
    \item Pero hay muchos colores que pueden ser considerados blancos.
    \item CIE estandariza los diferentes blancos.
    \item El blanco D65 es el que emitiría un cuerpo negro a 6500 K.
\end{itemize}
% \chapter{Iluminación y Sombreado}
% \chapter{Render}
% \chapter{Animación}
% \chapter{Visualización Científica}
% \chapter{Informática Gráfica en la Inteligencia Artificial}

\end{document}